\section{Where the numbers on the slides come from}
\label{app:provenance}

The slides were locked on 4 November 2026 and are submitted unchanged, as the
brief requires. At the presentation we were asked where the numbers on the
business-value slide (slide 11) come from; this appendix answers that for every
number on every slide. Table~\ref{tab:provenance} lists each one, slide 11 first,
with its kind, its source or formula and, where the slide's wording is wrong or
loose, the reading this report stands by. \emph{Observed} numbers are counted in
the Olist tables, on the test window (orders placed 30 June to 30 August 2018:
12{,}927 orders, 12{,}855 of them reviewed, 12{,}630 delivered) unless the row
or its slide's heading names another window. \emph{Assumed} numbers have no data source: Olist publishes
no costs, so they are inputs we chose, and Appendix~\ref{app:customer} shows what
changes if they are wrong. \emph{Derived} numbers are arithmetic on the other
two, with the formula given. \emph{Planned} numbers are design choices: a pilot that has
not been run, or the size of a recommended list. Three further kinds carry no number to check: a claim taken from the
\emph{literature}, a claim with \emph{no data source} that the tables can neither
support nor refute, and \emph{wording} whose reading needs correcting. Of the 105
rows over the 13 slides, 88 carry a number that \code{src/slide\_provenance.py}
recomputes from the repository's outputs into
\code{report/results/slide\_provenance.csv}. All 88 match the slide to half a
unit of its last digit, except three figures the slides rounded loosely
(100{,}000 orders, 75\% and the 28 of 27--28\%), which match within the
wider tolerance recorded in the CSV; the value after the arrow is the recomputed
one. The script stops if any no longer matches, and a validation check fails if
the CSV records a mismatch or a cost that is not labelled as assumed.

% Generated by src/slide_provenance.py from report/results/slide_provenance.csv.
% Do not edit by hand: change the script and run it again.
{\footnotesize
\setlength{\tabcolsep}{3pt}
\renewcommand{\arraystretch}{1.1}
\begin{longtable}{@{}p{\dimexpr(\textwidth-6\tabcolsep)*24/100\relax}p{\dimexpr(\textwidth-6\tabcolsep)*10/100\relax}p{\dimexpr(\textwidth-6\tabcolsep)*33/100\relax}p{\dimexpr(\textwidth-6\tabcolsep)*33/100\relax}@{}}
\caption{Where the numbers on the slides come from: every number or claim, slide 11 first, with its kind, its source or formula ($\rightarrow$ the value recomputed from it) and the reading this report stands by. Rows count the test window unless the slide's heading or a bracket names another. Generated by \code{src/slide\_provenance.py} from \code{report/results/slide\_provenance.csv}.}\label{tab:provenance}\tabularnewline
\toprule
\textbf{On the slide} & \textbf{Kind} & \textbf{Source or formula} & \textbf{Reading or correction}\tabularnewline
\midrule
\endfirsthead
\multicolumn{4}{@{}l}{\textit{Table~\ref{tab:provenance}, continued}}\tabularnewline
\toprule
\textbf{On the slide} & \textbf{Kind} & \textbf{Source or formula} & \textbf{Reading or correction}\tabularnewline
\midrule
\endhead
\bottomrule
\endfoot
\multicolumn{4}{@{}p{\textwidth}@{}}{\textbf{Slide 11}: business value and the pilot}\\*
\raggedright ``Worth R\$775–R\$1,466 /\allowbreak{} 1k orders'' & \raggedright derived & \raggedright review\_\allowbreak{}timing\_\allowbreak{}value.csv, threshold 0.11: net\_\allowbreak{}reachable\_\allowbreak{}per\_\allowbreak{}1k, net\_\allowbreak{}all\_\allowbreak{}per\_\allowbreak{}1k $\rightarrow$ 775.2; 1,466.0 & \raggedright per 1,000 reviewed test orders, not orders; ceilings that count every contact as turning its customer around (at one in four: $-$R\$368 and $-$R\$196, customer\_\allowbreak{}value.csv)\tabularnewline
\raggedright ``R\$5 CALL'' & \raggedright assumed & \raggedright config.COSTS['outreach\_\allowbreak{}cost'] $\rightarrow$ 5 & \raggedright no data source; varied R\$2.5–10 in customer\_\allowbreak{}value\_\allowbreak{}sensitivity.csv\tabularnewline
\raggedright ``R\$40 BAD REVIEW'' & \raggedright assumed & \raggedright config.COSTS['low\_\allowbreak{}review\_\allowbreak{}cost'] $\rightarrow$ 40 & \raggedright stands for a 1–2$\star$ rating's effect on later buyers; the customer's own return cannot price it (repeat within 180 days: 1.50\% after 1–2$\star$, 1.90\% after 4–5$\star$, customer\_\allowbreak{}repeat.csv); varied R\$20–80\tabularnewline
\raggedright ``R\$775 conservative ceiling on the 35/\allowbreak{}100 share'' & \raggedright derived & \raggedright 1000 $\times$ (490 $\times$ 40 $-$ 1,927 $\times$ 5) /\allowbreak{} 12,855: the threshold-0.11 list's 490 hits not written before arrival (282 delivered + 208 never delivered) $\rightarrow$ 775.2 & \raggedright belongs to the threshold-0.11 list (31 contacts a day), not to the 35-in-100 share, which is worth R\$433\tabularnewline
\raggedright ``R\$433 at fixed 10 calls/\allowbreak{}day'' & \raggedright derived & \raggedright 1000 $\times$ (216 $\times$ 40 $-$ 614 $\times$ 5) /\allowbreak{} 12,855 $\rightarrow$ 433.3 & \raggedright of the 35 in 100, 26.4 never received the parcel; only 8.8 could be reached at delivery (that call alone: $-$R\$71)\tabularnewline
\raggedright ``35/\allowbreak{}100 share'' & \raggedright observed & \raggedright review\_\allowbreak{}timing\_\allowbreak{}capacity.csv, 10 a day: 216 of 614 = 162 never delivered + 54 delivered, not yet written $\rightarrow$ 35.2 & \raggedright includes the orders that never arrived\tabularnewline
\raggedright ``R\$1,466 full ceiling across all hits'' & \raggedright derived & \raggedright cost\_\allowbreak{}p2.csv, threshold 0.11: 1000 $\times$ (712 $\times$ 40 $-$ 1,927 $\times$ 5) /\allowbreak{} 12,855 $\rightarrow$ 1,466.0 & \raggedright 0.11 is the best of 91 thresholds read on the test window; the threshold fixed on validation, 0.19, gives R\$1,373 (R\$692 not written); 60.4\% of the gross value is never-delivered orders or late parcels rated before arrival\tabularnewline
\raggedright ``50 /\allowbreak{} 50 for 6 weeks (191 orders per arm)'' & \raggedright planned & \raggedright pilot\_\allowbreak{}design.csv (k 10, 25\%): weeks = ceil(2 $\times$ 191 /\allowbreak{} 10 /\allowbreak{} 7); n\_\allowbreak{}per\_\allowbreak{}arm\_\allowbreak{}80 $\rightarrow$ 6; 191 & \raggedright ten a day on every calendar day (8 weeks on weekdays); re-sized per contact in customer\_\allowbreak{}pilot\_\allowbreak{}arms.csv: message 209 per arm, 5 weeks; delivery-day call 159, 5 weeks\tabularnewline
\raggedright ``5 RANDOMLY CALLED /\allowbreak{} DAY'' \ldots{} ``5 HELD AS CONTROL /\allowbreak{} DAY'' & \raggedright planned & \raggedright 1:1 randomisation inside the daily top ten $\rightarrow$ 5; 5 & \raggedright as on the slide\tabularnewline
\raggedright ``TODAY (NO CALL) /\allowbreak{} 57.2\%'' & \raggedright observed & \raggedright capacity\_\allowbreak{}p2\_\allowbreak{}daily.csv (k 10): 351 of 614 $\rightarrow$ 57.17 & \raggedright a historical rate, not a control group (no pilot ran): 26.4 points never received the parcel and 22.0 had rated a late parcel while waiting\tabularnewline
\raggedright ``BREAK-EVEN ($\sim$1 IN 5) /\allowbreak{} $\leq$ 44.7\%'' & \raggedright derived & \raggedright pilot\_\allowbreak{}breakeven.csv: 0.572 /\allowbreak{} 0.125; 0.572 $-$ 5/\allowbreak{}40 $\rightarrow$ 4.6; 44.67 & \raggedright 21.9\% (1 in 4.6); unrelated to the 44.7\% of customers the 11\% late warning reaches\tabularnewline
\raggedright ``6-WK TARGET (1 IN 4) /\allowbreak{} 42.9\% /\allowbreak{} 80\% power'' & \raggedright planned & \raggedright pilot\_\allowbreak{}design.csv: 0.572 $\times$ (1 $-$ 0.25); power 0.80, alpha 0.05 $\rightarrow$ 42.87; 80 & \raggedright one in four is a design choice, not an estimate\tabularnewline
\raggedright ``+R\$0.72/\allowbreak{}call'' & \raggedright derived & \raggedright (0.572 $-$ 0.429) $\times$ 40 $-$ 5 $\rightarrow$ 0.717 & \raggedright only if the one-in-four effect and both costs hold; the ten-a-day list then nets R\$34 per 1,000 reviewed orders\tabularnewline
\raggedright ``Nobody in the 2016–2018 log'' & \raggedright no data source & \raggedright none & \raggedright the slide's statement about calls has no source: the tables record no contact other than the survey, so only a randomised pilot can measure what a contact changes\tabularnewline
\raggedright ``PLUS ARM B (PROMISED-DATE OVERDUE ORDERS)'' & \raggedright planned & \raggedright customer\_\allowbreak{}pilot\_\allowbreak{}arms.csv & \raggedright 13.0 customers a day rate the order only after the promised day ends without a parcel, 54.7\% of them 1–2$\star$: 209 per arm, 5 weeks (two arms) or 7 (three); outcome: the final rating\tabularnewline
\raggedright ``LIVE VALIDATION'' \ldots{} ``Historical data proves'' & \raggedright wording & \raggedright none & \raggedright nothing was validated live: the held-out days show; the pilot would measure\tabularnewline
\addlinespace[3pt]
\multicolumn{4}{@{}p{\textwidth}@{}}{\textbf{Slide 1}: title}\\*
\raggedright ``100,000 Brazilian e-commerce orders'' \ldots{} ``99,441 orders (2016–2018)'' & \raggedright observed & \raggedright orders table: rows, purchase years $\rightarrow$ 99,441; 99,441; 2016; 2018 [whole extract, 2016–18] & \raggedright the analysis uses the 99,091 orders placed January 2017 to August 2018\tabularnewline
\raggedright ``3. Delivery /\allowbreak{} Call top-10 risk list'' & \raggedright planned & \raggedright customer\_\allowbreak{}contact\_\allowbreak{}lists.csv (delivery\_\allowbreak{}day\_\allowbreak{}k10) $\rightarrow$ 10 & \raggedright worked on the previous day's deliveries not yet rated: 26.0 in 100 rate 1–2$\star$ (3.2$\times$ the 8.1\% base)\tabularnewline
\addlinespace[3pt]
\multicolumn{4}{@{}p{\textwidth}@{}}{\textbf{Slide 2}: one customer's order \textit{(rows: order e1320973, unless bracketed)}}\\*
\raggedright ``16 Jul 2018 /\allowbreak{} Day 0'' \ldots{} ``7 Aug 2018 /\allowbreak{} Day 22'' \ldots{} ``9 Aug 2018 /\allowbreak{} Day 24'' \ldots{} ``11 Aug 2018 /\allowbreak{} Day 26'' \ldots{} ``17 Aug 2018 /\allowbreak{} Day 32'' & \raggedright observed & \raggedright customer\_\allowbreak{}story.csv: purchase, promised, survey\_\allowbreak{}sent, review\_\allowbreak{}answered, delivered $\rightarrow$ 16 Jul 2018; 0; 7 Aug 2018; 22; 9 Aug 2018; 24; 11 Aug 2018; 26; 17 Aug 2018; 32 & \raggedright as on the slide\tabularnewline
\raggedright ``Order \#e1320973  ·  Health \& beauty item (R\$65 + R\$28 freight)'' & \raggedright observed & \raggedright customer\_\allowbreak{}story.csv: order\_\allowbreak{}id, category, price, freight $\rightarrow$ e1320973; health \& beauty; 65.0; 27.7 & \raggedright R\$64.99 + R\$27.74\tabularnewline
\raggedright ``2,267 km from São Luís (MA) to São Gonçalo (RJ)'' & \raggedright observed & \raggedright customer\_\allowbreak{}story.csv: distance\_\allowbreak{}km (postcode-prefix centroids), seller and customer city $\rightarrow$ 2,266.9; sao luis (MA); sao goncalo (RJ) & \raggedright as on the slide\tabularnewline
\raggedright ``($\sim$21-day window)'' \ldots{} ``(10.5 days late)'' \ldots{} ``5.8 days old'' \ldots{} ``took 31.8 days'' & \raggedright observed & \raggedright customer\_\allowbreak{}story.csv: promised\_\allowbreak{}days, days\_\allowbreak{}late, review\_\allowbreak{}days\_\allowbreak{}before\_\allowbreak{}arrival, actual\_\allowbreak{}days $\rightarrow$ 21.3; 10.50; 5.77; 31.79 & \raggedright as on the slide\tabularnewline
\raggedright ``Zero proactive warning sent to the buyer.'' & \raggedright no data source & \raggedright none & \raggedright the tables record no message and no contact other than the survey\tabularnewline
\raggedright ``Triggered by the promised calendar date'' & \raggedright no data source & \raggedright customer\_\allowbreak{}story.csv; customer\_\allowbreak{}survey\_\allowbreak{}timing.csv & \raggedright the trigger is not recorded; the survey went out 2 days after the promised date, 8.5 days before delivery (test window: 71.4\% of late 1–2$\star$ parcels exactly 2 days after)\tabularnewline
\raggedright ``Frustrated and still empty-handed'' & \raggedright wording & \raggedright customer\_\allowbreak{}story.csv & \raggedright rests on the timestamps (answered 11 Aug 17:33, delivered 17 Aug 11:55); the comment (`Bom' /\allowbreak{} `Não veio a embalagem certa') is ambiguous and may mean the packaging\tabularnewline
\raggedright ``the damage on the product page is already done'' & \raggedright no data source & \raggedright none & \raggedright the tables do not record what later buyers saw\tabularnewline
\addlinespace[3pt]
\multicolumn{4}{@{}p{\textwidth}@{}}{\textbf{Slide 3}: broken promises and ratings \textit{(rows: 2017–18, delivered and reviewed, unless bracketed)}}\\*
\raggedright ``ACROSS 95,561 DELIVERED ORDERS'' & \raggedright observed & \raggedright review\_\allowbreak{}timing.csv (all): n $\rightarrow$ 95,561 & \raggedright delivered and reviewed orders\tabularnewline
\raggedright ``DELIVERED ON TIME /\allowbreak{} 9.2\%'' \ldots{} ``87,902 orders (92.0\%)'' & \raggedright observed & \raggedright review\_\allowbreak{}timing.csv (on time): pct\_\allowbreak{}low\_\allowbreak{}review, n, share $\rightarrow$ 9.20; 87,902; 91.99 & \raggedright as on the slide\tabularnewline
\raggedright ``DELIVERED LATE /\allowbreak{} 54.0\%'' \ldots{} ``7,659 orders (8.0\%)'' & \raggedright observed & \raggedright review\_\allowbreak{}timing.csv (late): pct\_\allowbreak{}low\_\allowbreak{}review, n, share $\rightarrow$ 54.00; 7,659; 8.01 & \raggedright as on the slide\tabularnewline
\raggedright ``5.9$\times$'' & \raggedright derived & \raggedright 54.0 /\allowbreak{} 9.2 $\rightarrow$ 5.87 & \raggedright as on the slide\tabularnewline
\raggedright ``78.6\%'' \ldots{} ``78.0\%'' \ldots{} ``61.3\%'' \ldots{} ``19.1\%'' \ldots{} ``9.7\%'' \ldots{} ``8.9\%'' & \raggedright observed & \raggedright review\_\allowbreak{}timing\_\allowbreak{}lateness.csv: low\_\allowbreak{}pct by days late or early $\rightarrow$ 78.63; 77.99; 61.31; 19.13; 9.68; 8.93 & \raggedright as on the slide\tabularnewline
\raggedright ``1–3 days late doubles bad reviews (19.1\%); 4+ days late jumps past 61\%–79\%'' & \raggedright derived & \raggedright 19.1 /\allowbreak{} 9.2; review\_\allowbreak{}timing\_\allowbreak{}lateness.csv $\rightarrow$ 2.1; 19.13; 61.3; 78.6 & \raggedright by day: 12.3\% during the promised day, 19.8\% one to two days late (customer\_\allowbreak{}lateness\_\allowbreak{}days.csv)\tabularnewline
\raggedright ``Only 1 in 12 parcels (8.0\%) arrives late'' \ldots{} ``34\% (one-third) of every 1–2$\star$ review'' & \raggedright observed & \raggedright review\_\allowbreak{}timing.csv: late share; share\_\allowbreak{}of\_\allowbreak{}all\_\allowbreak{}low\_\allowbreak{}reviews\_\allowbreak{}pct $\rightarrow$ 12.5; 8.01; 33.8 & \raggedright of the 1–2$\star$ reviews on delivered orders\tabularnewline
\raggedright ``On Olist's own data, bad logistics reviews directly depress a seller's future sales'' & \raggedright literature & \raggedright refs.bib deshpande2023logistics & \raggedright the study is of another marketplace, not Olist; seller sales were not measured here\tabularnewline
\addlinespace[3pt]
\multicolumn{4}{@{}p{\textwidth}@{}}{\textbf{Slide 4}: when the review is written \textit{(rows: 2017–18, delivered and reviewed, unless bracketed)}}\\*
\raggedright ``ON-TIME PARCELS  (11 IN 12)'' \ldots{} ``LATE PARCELS  (1 IN 12)'' & \raggedright observed & \raggedright review\_\allowbreak{}timing.csv: 12 $\times$ each share $\rightarrow$ 11.0; 1.0 & \raggedright as on the slide\tabularnewline
\raggedright ``+2.1 days /\allowbreak{} after delivery scan'' & \raggedright observed & \raggedright review\_\allowbreak{}timing.csv (low review, on time): median days from delivery to the review $\rightarrow$ 2.12 & \raggedright on-time parcels rated 1–2$\star$\tabularnewline
\raggedright ``66\% of all 1–2$\star$ reviews /\allowbreak{} Wrong, damaged, or disappointing items'' \ldots{} ``99.5\% review AFTER delivery'' & \raggedright observed & \raggedright review\_\allowbreak{}timing.csv (low review, on time): share of 1–2$\star$, written after arrival $\rightarrow$ 66.2; 99.51 & \raggedright not measured on 4 November; since: of on-time 1–2$\star$ comments 34.0\% mention a wrong or missing item, 15.0\% damage or a defect, 16.9\% quality (keyword topics, customer\_\allowbreak{}complaint\_\allowbreak{}topics.csv)\tabularnewline
\raggedright ``1$\star$ written 4.5d early!'' & \raggedright observed & \raggedright review\_\allowbreak{}timing.csv (low review, late): median days from the review to delivery $\rightarrow$ 4.50 & \raggedright all late 1–2$\star$, not only 1$\star$\tabularnewline
\raggedright ``ALREADY WRITTEN /\allowbreak{} 87.3\%'' & \raggedright observed & \raggedright review\_\allowbreak{}timing.csv (low review, late): pct\_\allowbreak{}written\_\allowbreak{}before\_\allowbreak{}arrival $\rightarrow$ 87.31 & \raggedright 77.4\% on the test window\tabularnewline
\raggedright ``34\% of all 1–2$\star$ reviews /\allowbreak{} Survey fires at the promised date'' & \raggedright observed & \raggedright review\_\allowbreak{}timing.csv (low review, late); customer\_\allowbreak{}survey\_\allowbreak{}timing.csv $\rightarrow$ 33.8 & \raggedright survey on or after the promised date for 99.5\%, a median 2 days after; the rule is not recorded\tabularnewline
\addlinespace[3pt]
\multicolumn{4}{@{}p{\textwidth}@{}}{\textbf{Slide 5}: why promises break \textit{(rows: monthly, 2017–18, unless bracketed)}}\\*
\raggedright ``sped up 2$\times$'' \ldots{} ``16.9 d $\rightarrow$ 7.7 d average delivery'' & \raggedright observed & \raggedright lateness\_\allowbreak{}monthly.csv mean\_\allowbreak{}lead, Feb and Aug 2018 $\rightarrow$ 2.2; 16.95; 7.73 & \raggedright delivery time halved\tabularnewline
\raggedright ``75\% of total time lives with the carrier'' & \raggedright observed & \raggedright orders table: mean (delivered $-$ carrier handover) /\allowbreak{} mean lead time $\rightarrow$ 74.5 [2017–18, delivered] & \raggedright 74.5\% (9.3 of 12.5 days), rounded up\tabularnewline
\raggedright ``across 20 months (Jan 2017 – Aug 2018)'' & \raggedright observed & \raggedright lateness\_\allowbreak{}monthly.csv rows $\rightarrow$ 20 & \raggedright as on the slide\tabularnewline
\raggedright ``Carrier jam peak 21.4\%'' \ldots{} ``Black Friday 14.3\%'' \ldots{} ``16.0\%'' \ldots{} ``10.4\%'' & \raggedright observed & \raggedright lateness\_\allowbreak{}monthly.csv late\_\allowbreak{}pct: Mar 2018, Nov 2017, Feb 2018, Aug 2018 $\rightarrow$ 21.36; 14.31; 16.00; 10.39 & \raggedright as on the slide\tabularnewline
\raggedright ``$\sim$11 extra buffer days on average'' & \raggedright observed & \raggedright orders table: mean promised days $-$ mean lead time $\rightarrow$ 11.1 [2017–18, delivered] & \raggedright 23.6 $-$ 12.5; test window 9.1 (17.48 $-$ 8.34)\tabularnewline
\raggedright ``(9.8d error)'' & \raggedright observed & \raggedright reg\_\allowbreak{}p1a.csv: Olist's own delivery promise, MAE $\rightarrow$ 9.79 [test window] & \raggedright a test-window figure; the $\sim$11 days are 2017–18\tabularnewline
\raggedright ``missed 1 in 5 orders (21.4\%) when regional carrier slowdowns struck'' & \raggedright observed & \raggedright lateness\_\allowbreak{}monthly.csv (Mar 2018); lateness\_\allowbreak{}peak\_\allowbreak{}by\_\allowbreak{}region.csv $\rightarrow$ 4.7; 21.36 & \raggedright network-wide: every region 17.0–30.5\% late in Feb–Mar 2018 (5.3–11.8\% otherwise)\tabularnewline
\raggedright ``inflates validation AUC from 0.59 to a fake 1.00'' & \raggedright observed & \raggedright report/\allowbreak{}stats.txt: c1b\_\allowbreak{}logreg\_\allowbreak{}roc, c1b\_\allowbreak{}leaked\_\allowbreak{}roc $\rightarrow$ 0.589; 1.000 [test window] & \raggedright a test-window score of the Phase 1 logistic regression, not a validation AUC\tabularnewline
\raggedright ``unseen Jul–Aug 2018 orders'' & \raggedright wording & \raggedright config.VAL\_\allowbreak{}END, config.ANALYSIS\_\allowbreak{}END & \raggedright orders placed 30 June to 30 August 2018: 12,927 (12,855 reviewed, 12,630 delivered)\tabularnewline
\raggedright ``(19.5 $\rightarrow$ 36.8 days), hiding June lateness (1.4\%)'' & \raggedright observed & \raggedright lateness\_\allowbreak{}weekly.csv promise\_\allowbreak{}days, weeks of 14 and 28 May 2018; lateness\_\allowbreak{}monthly.csv, June 2018 $\rightarrow$ 19.49; 36.80; 1.36 [May–June 2018] & \raggedright weekly means; `hiding' is a reading\tabularnewline
\addlinespace[3pt]
\multicolumn{4}{@{}p{\textwidth}@{}}{\textbf{Slide 6}: three moments, three tools \textit{(rows: 2017–18, delivered and reviewed, unless bracketed)}}\\*
\raggedright ``12.8\% /\allowbreak{} ($\sim$1 in 8 delivered orders)'' \ldots{} ``14.6\% across all reviewed orders including undelivered'' & \raggedright observed & \raggedright review\_\allowbreak{}timing.csv (all): pct\_\allowbreak{}low\_\allowbreak{}review; 14,365 /\allowbreak{} 98,329 reviewed (customer\_\allowbreak{}complaint\_\allowbreak{}topics.csv, train\_\allowbreak{}val\_\allowbreak{}test.csv) $\rightarrow$ 12.79; 7.8; 14.61 [2017–18, reviewed] & \raggedright as on the slide\tabularnewline
\raggedright ``a single decision tree captures 96\% of the delivery-triage signal'' & \raggedright derived & \raggedright clf\_\allowbreak{}p2.csv PR-AUC: tree /\allowbreak{} deployed stack $\rightarrow$ 95.7 [test window] & \raggedright review risk, not delivery: 0.396 against 0.414\tabularnewline
\raggedright ``Replaces padded 24-day guesses'' & \raggedright observed & \raggedright orders table: mean promised days $\rightarrow$ 23.6 [2017–18, delivered] & \raggedright 23.6 days (2017–18), 17.5 on test; the recommended date is the forecast + 8 days (18.6 on average)\tabularnewline
\raggedright ``pauses the review email before 87.3\% of late 1–2$\star$ are written'' & \raggedright observed & \raggedright as slide 4 $\rightarrow$ 87.31 & \raggedright recommended: ask for the review once the order is complete; every customer is still asked; no review is filtered, delayed in publication or rewarded\tabularnewline
\raggedright ``Triggers automatically the instant a promised date passes'' & \raggedright wording & \raggedright customer\_\allowbreak{}outcomes\_\allowbreak{}test.csv; customer\_\allowbreak{}lateness\_\allowbreak{}days.csv & \raggedright the trigger waits for the end of the promised day: 335 of 936 late parcels arrived during it, rated like on-time ones (12.3\% 1–2$\star$, against 11.1\% a day earlier)\tabularnewline
\raggedright ``$\sim$2.1-day post-delivery window'' & \raggedright observed & \raggedright as slide 4 $\rightarrow$ 2.12 & \raggedright as on the slide\tabularnewline
\addlinespace[3pt]
\multicolumn{4}{@{}p{\textwidth}@{}}{\textbf{Slide 7}: results at checkout}\\*
\raggedright ``Tested on 12,630 unseen future deliveries'' & \raggedright observed & \raggedright predictions\_\allowbreak{}p1a\_\allowbreak{}test.csv rows $\rightarrow$ 12,630 & \raggedright delivered test orders\tabularnewline
\raggedright ``cuts date error by 60\%'' \ldots{} ``9.79 days'' \ldots{} ``3.90 days  (beats GBDT 4.14d \& Linear 4.18d)'' & \raggedright derived & \raggedright reg\_\allowbreak{}p1a.csv MAE; 1 $-$ 3.90 /\allowbreak{} 9.79 $\rightarrow$ 60.2; 9.788; 3.898; 4.141; 4.181 & \raggedright as on the slide\tabularnewline
\raggedright ``Costs R\$0 to quote'' & \raggedright assumed & \raggedright none & \raggedright the forecast on its own is exceeded by 31.9\% of deliveries (the platform's date by 7.4\%); forecast + 8 days: 2.5\%, 1.1 days longer (customer\_\allowbreak{}promise\_\allowbreak{}rules.csv)\tabularnewline
\raggedright ``Pads $\sim$11 extra buffer days'' & \raggedright observed & \raggedright as slide 5 $\rightarrow$ 11.1 [2017–18, delivered] & \raggedright 9.1 on the test window\tabularnewline
\raggedright ``lifts PR-AUC 2.9$\times$ vs 1.2$\times$ (ROC-AUC 0.782 vs 0.584)'' & \raggedright observed & \raggedright clf\_\allowbreak{}p1b.csv: derived route vs HistGradientBoostingClassifier $\rightarrow$ 2.89; 1.24; 0.7820; 0.5838 & \raggedright as on the slide\tabularnewline
\raggedright ``35\% of late orders (3.5$\times$)'' & \raggedright observed & \raggedright capacity\_\allowbreak{}p1b.csv (derived, top 10\%): recall, lift $\rightarrow$ 34.8; 3.48 & \raggedright 25.8\% of that tenth were late\tabularnewline
\raggedright ``\#23 of 238 delivered orders on 16 July'' & \raggedright observed & \raggedright customer\_\allowbreak{}story.csv (p\_\allowbreak{}late\_\allowbreak{}derived rank) $\rightarrow$ 23; 238 [16 July 2018 purchases] & \raggedright as on the slide\tabularnewline
\raggedright ``$\geq$ 11\% risk: Pre-warns 83\% (5 in 6)'' & \raggedright observed & \raggedright capacity\_\allowbreak{}operating\_\allowbreak{}points.csv (P1B, 0.11): recall $\rightarrow$ 11; 82.6; 5.0 & \raggedright reaches 44.7\% of customers; 86.3\% of them then on time (773 late, 4,874 on time)\tabularnewline
\raggedright ``R\$12 express upgrade vs. R\$40 bad-review cost'' & \raggedright assumed & \raggedright config.COSTS['expedite\_\allowbreak{}cost'], ['late\_\allowbreak{}delivery\_\allowbreak{}cost'] $\rightarrow$ 12; 40 & \raggedright R\$40 is charged for every late parcel\tabularnewline
\raggedright ``27–28\% precision  $<$  30\% break-even'' & \raggedright observed & \raggedright capacity\_\allowbreak{}p1b.csv precision, top 1–7.5\%; 12 /\allowbreak{} 40 $\rightarrow$ 26.8; 27.4; 30 & \raggedright 26.8–27.4\% (25.8\% at the top 10\%)\tabularnewline
\raggedright ``at least 30\% (12 /\allowbreak{} 40)'' & \raggedright derived & \raggedright 12 /\allowbreak{} 40 $\rightarrow$ 30 & \raggedright conservative: with the observed 1–2$\star$ gap (54.0\% $-$ 9.2\%) a late parcel costs R\$17.9, break-even 0.67\tabularnewline
\raggedright ``75\% of transit happens after checkout inside the carrier network'' & \raggedright observed & \raggedright as slide 5 $\rightarrow$ 74.5 [2017–18, delivered] & \raggedright 74.5\%, after the seller's handover\tabularnewline
\addlinespace[3pt]
\multicolumn{4}{@{}p{\textwidth}@{}}{\textbf{Slide 8}: results at delivery}\\*
\raggedright ``Across 62 unseen test days (12,855 reviewed orders)'' & \raggedright observed & \raggedright predictions\_\allowbreak{}p2\_\allowbreak{}test.csv: purchase days, rows $\rightarrow$ 62; 12,855 & \raggedright as on the slide\tabularnewline
\raggedright ``turns 11\% random hits into 57\%'' \ldots{} ``(57.2\%)'' \ldots{} ``(11.0\% base)'' \ldots{} ``89 of 100 calls wasted on happy buyers'' & \raggedright observed & \raggedright capacity\_\allowbreak{}p2\_\allowbreak{}daily.csv (k 10); review\_\allowbreak{}timing\_\allowbreak{}capacity.csv base\_\allowbreak{}rate $\rightarrow$ 11.0; 57.2; 57.17; 11.03; 89.0 & \raggedright a historical outcome (no contact recorded besides the survey) of a purchase-day list scored with the realised delivery; deliveries not yet rated: 26.0 per 100; 3–5$\star$ is not all `happy'\tabularnewline
\raggedright ``5.2$\times$ at 10 calls/\allowbreak{}day'' & \raggedright derived & \raggedright 0.572 /\allowbreak{} 0.110 $\rightarrow$ 5.18 & \raggedright as on the slide\tabularnewline
\raggedright ``(68\% at top 5\%)'' & \raggedright observed & \raggedright capacity\_\allowbreak{}p2.csv (top 5\%): precision $\rightarrow$ 67.8 & \raggedright 43.1\% of those 643 never arrived; 31.6 in 100 had rated before the parcel arrived; 4.8 in 100 could be reached at delivery\tabularnewline
\raggedright ``Powered by delay vs. promise, multi-unit baskets, high freight share \& category history.'' & \raggedright observed & \raggedright perm\_\allowbreak{}p2.csv & \raggedright delay vs promise 0.206, days 0.074, items 0.045; freight share and category history rank 32 and 35 of 36\tabularnewline
\raggedright ``still hits 45.1\% unhappy buyers (vs 9.6\% at random)'' & \raggedright observed & \raggedright review\_\allowbreak{}timing\_\allowbreak{}composition.csv (delivered only) $\rightarrow$ 45.08; 9.63 & \raggedright 28.9 of the 45.1 had rated before their late parcel arrived; 16.2 per 100 not yet rated at delivery\tabularnewline
\raggedright ``top 5\% = 28\% (3.5$\times$), where 99.5\% wait $\sim$2.1 days after delivery'' & \raggedright observed & \raggedright review\_\allowbreak{}timing\_\allowbreak{}subsets.csv (on time); review\_\allowbreak{}timing.csv $\rightarrow$ 28.0; 3.53; 99.51; 2.12 & \raggedright 99.5\% and 2.1 days are 2017–18 figures\tabularnewline
\raggedright ``Only 13 hits /\allowbreak{} 100 needed'' & \raggedright derived & \raggedright 100 $\times$ 5 /\allowbreak{} 40 $\rightarrow$ 12.5 & \raggedright 12.5 if every contact works; 12.5 /\allowbreak{} s if a share s does\tabularnewline
\raggedright ``57 total hits (and 35 non-pre-empted) clear the 12.5\% bar'' & \raggedright observed & \raggedright review\_\allowbreak{}timing\_\allowbreak{}capacity.csv (10 a day) $\rightarrow$ 57.2; 35.2; 12.50 & \raggedright 26.4 of the 35 never received the parcel; a call at delivery reaches 8.8, below 12.5\tabularnewline
\raggedright ``35 not written before arrival'' \ldots{} ``22 wrote waiting'' \ldots{} ``43 satisfied (3–5$\star$)'' \ldots{} ``614 total calls over 62 days'' & \raggedright observed & \raggedright review\_\allowbreak{}timing\_\allowbreak{}capacity.csv (10 a day) $\rightarrow$ 35.2; 22.0; 42.8; 614; 62 & \raggedright list entries, not calls; 3–5$\star$ is not all `satisfied'\tabularnewline
\addlinespace[3pt]
\multicolumn{4}{@{}p{\textwidth}@{}}{\textbf{Slide 9}: the app and 16 July 2018 \textit{(rows: 16 July 2018 purchases, unless bracketed)}}\\*
\raggedright ``8-tab Streamlit app'' \ldots{} ``replay any of 62 test days'' & \raggedright observed & \raggedright predictions\_\allowbreak{}p2\_\allowbreak{}test.csv purchase days $\rightarrow$ 62 [test window] & \raggedright eight tabs when the slides were locked on 4 November 2026\tabularnewline
\raggedright ``16 JULY 2018  (244 REVIEWED ORDERS)'' \ldots{} ``8 of 10  real 1–2$\star$'' \ldots{} ``(vs $\sim$1 by random)'' & \raggedright observed & \raggedright customer\_\allowbreak{}story.csv; 10 $\times$ the day's 1–2$\star$ rate $\rightarrow$ 244; 8; 10; 1.3 & \raggedright a purchase-day list scored with the realised delivery: all 8 reachable only when the promised day ended (4 never delivered, 4 rated while waiting); 0 at delivery\tabularnewline
\raggedright ``Checkout error: 4.3 d vs Olist's 10.7 d'' & \raggedright observed & \raggedright predictions\_\allowbreak{}p1a\_\allowbreak{}test.csv, 16 July: MAE of forecast and promise $\rightarrow$ 4.35; 10.72 & \raggedright the day's 238 delivered orders\tabularnewline
\raggedright ``4 buyers wrote 1$\star$ while still waiting for a late parcel'' & \raggedright observed & \raggedright customer\_\allowbreak{}journeys\_\allowbreak{}test.csv, 16 July top ten $\rightarrow$ 4 & \raggedright as on the slide\tabularnewline
\raggedright ``4 buyers never received their parcel at all (shipped, invoiced or cancelled)'' & \raggedright observed & \raggedright customer\_\allowbreak{}journeys\_\allowbreak{}test.csv, 16 July top ten: order\_\allowbreak{}status $\rightarrow$ 4 & \raggedright 2 shipped, 1 invoiced, 1 cancelled\tabularnewline
\raggedright ``57 /\allowbreak{} 100 average precision'' & \raggedright observed & \raggedright capacity\_\allowbreak{}p2\_\allowbreak{}daily.csv (k 10) $\rightarrow$ 57.2 [test window] & \raggedright 351 of 614, pooled\tabularnewline
\raggedright ``(85\% RISK)'' \ldots{} ``Flagged \#23 of 238 delivered at checkout, and \#1 of 244 reviewed on the call list'' & \raggedright observed & \raggedright customer\_\allowbreak{}story.csv $\rightarrow$ 84.7; 23; 238; 1; 244 & \raggedright `that morning's list' was scored once the delivery was known (17 August); at checkout its late risk was 18.2\%, above 11\%\tabularnewline
\addlinespace[3pt]
\multicolumn{4}{@{}p{\textwidth}@{}}{\textbf{Slide 10}: the recommendations}\\*
\raggedright ``Replace padded promises with our checkout forecast (60\% smaller error: 9.79 $\rightarrow$ 3.90 d)'' & \raggedright derived & \raggedright reg\_\allowbreak{}p1a.csv MAE $\rightarrow$ 60.2; 9.788; 3.898 & \raggedright the forecast on its own is exceeded by 31.9\% of deliveries (the platform's date by 7.4\%): recommended, forecast + 8 days (2.5\%)\tabularnewline
\raggedright ``banner above 11\% risk — pre-warning 83\% (5 in 6)'' & \raggedright observed & \raggedright as slide 7 $\rightarrow$ 11; 82.6; 5.0 & \raggedright reaches 44.7\%, 86.3\% then on time; not needed with the forecast-plus-margin date\tabularnewline
\raggedright ``closing the window behind 87.3\% of late 1–2$\star$'' & \raggedright observed & \raggedright as slide 4 $\rightarrow$ 87.31 [2017–18, delivered and reviewed] & \raggedright the effect on ratings is for the pilot's third arm to measure; every customer is still asked; no review is filtered, delayed in publication or rewarded\tabularnewline
\raggedright ``Cost: $\sim$R\$0   ·   Immediate trust gain'' & \raggedright assumed & \raggedright none & \raggedright an assumed cost; the trust gain is a hypothesis\tabularnewline
\raggedright ``Stops pre-arrival 1$\star$'' & \raggedright wording & \raggedright none & \raggedright by construction; the effect on the final rating is unmeasured\tabularnewline
\raggedright ``top $\sim$10 orders each morning (57\% precision = 5.2$\times$ random; 45\% on delivered parcels alone)'' \ldots{} ``still wait $\sim$2.1 days'' & \raggedright observed & \raggedright as slides 4 and 8 $\rightarrow$ 10; 57.2; 5.18; 45.1; 2.12 & \raggedright as a team would work it: 26.0 per 100 (3.2$\times$ the 8.1\% base)\tabularnewline
\raggedright ``Checkout-only precision plateaus at 27–28\% across top-K tiers — below the 30\% break-even'' & \raggedright observed & \raggedright as slide 7 $\rightarrow$ 26.8; 27.4; 30 & \raggedright 26.8–27.4\%; 30\% is conservative\tabularnewline
\raggedright ``Cost: $\sim$R\$5/\allowbreak{}call'' & \raggedright assumed & \raggedright config.COSTS['outreach\_\allowbreak{}cost'] $\rightarrow$ 5 & \raggedright as on the slide\tabularnewline
\raggedright ``Refit monthly on a rolling 12-month window (because transit speed halved in 6 months)'' & \raggedright wording & \raggedright lateness\_\allowbreak{}monthly.csv mean\_\allowbreak{}lead & \raggedright the 12-month window was not tested (expanding windows were); delivery time halved (16.9 $\rightarrow$ 7.7 days)\tabularnewline
\addlinespace[3pt]
\multicolumn{4}{@{}p{\textwidth}@{}}{\textbf{Slide 12}: lessons}\\*
\raggedright ``squeezing 0.001 AUC from a complex ensemble'' & \raggedright observed & \raggedright clf\_\allowbreak{}p2.csv PR-AUC: boosted trees $-$ stack $\rightarrow$ 0.0011 & \raggedright PR-AUC; the deployed stack is 0.0011 below\tabularnewline
\raggedright ``a fake 1.00 AUC'' & \raggedright observed & \raggedright as slide 5 $\rightarrow$ 1.000 & \raggedright as on the slide\tabularnewline
\raggedright ``87.3\% of late-order bad reviews were written before the parcel even arrived'' & \raggedright observed & \raggedright as slide 4 $\rightarrow$ 87.31 [2017–18, delivered and reviewed] & \raggedright as on the slide\tabularnewline
\raggedright ``27–28\% ($<$ 30\% rush break-even)'' & \raggedright observed & \raggedright as slide 7 $\rightarrow$ 26.8; 27.4; 30 & \raggedright as on the slide\tabularnewline
\raggedright ``how to run the 6-week pilot'' & \raggedright planned & \raggedright pilot\_\allowbreak{}design.csv (k 10, 25\%) $\rightarrow$ 6 & \raggedright per contact: 5 weeks (message, two arms; 7 with three), 5 (delivery-day call)\tabularnewline
\raggedright ``75\% of delivery time happens inside third-party carriers'' & \raggedright observed & \raggedright as slide 5 $\rightarrow$ 74.5 [2017–18, delivered] & \raggedright 74.5\%\tabularnewline
\raggedright ``Northeast /\allowbreak{} North routes face structural distance disadvantages.'' & \raggedright observed & \raggedright segments\_\allowbreak{}p1a\_\allowbreak{}region.csv & \raggedright mean delivery 15.6 (North), 13.4 (Northeast), 7.2 days (Southeast); forecast error 6.8, 6.3, 3.4\tabularnewline
\addlinespace[3pt]
\multicolumn{4}{@{}p{\textwidth}@{}}{\textbf{Slide 13}: closing and demo}\\*
\raggedright ``60\% closer date'' \ldots{} ``Top-10 calls (5.2$\times$)'' & \raggedright derived & \raggedright as slides 7 and 8 $\rightarrow$ 60.2; 5.18 & \raggedright as on the slide\tabularnewline
\raggedright ``Saved before 1$\star$!'' & \raggedright wording & \raggedright none & \raggedright means reaching the customer before they rate the order; every customer is still asked; no review is filtered, delayed in publication or rewarded\tabularnewline
\raggedright ``8 of 10 real 1–2$\star$ (\#1 = São Gonçalo!)'' & \raggedright observed & \raggedright as slide 9 $\rightarrow$ 8; 10; 1 [16 July 2018 purchases] & \raggedright as on the slide\tabularnewline
\raggedright ``What-if cost \& pilot planner'' \ldots{} ``Test R\$5 call vs R\$40 review \& 6-wk RCT'' & \raggedright assumed & \raggedright config.COSTS; pilot\_\allowbreak{}design.csv weeks\_\allowbreak{}80 $\rightarrow$ 5; 40; 6 & \raggedright both costs are assumptions (customer\_\allowbreak{}value\_\allowbreak{}sensitivity.csv varies them); pilot sizes per contact: customer\_\allowbreak{}pilot\_\allowbreak{}arms.csv\tabularnewline
\end{longtable}}


\paragraph{The deck's vocabulary.} Some of the deck's words are not the
report's, and some say more than the data show. The adjective the slides attach
to the team's dates, test, tools, findings and caveats (slides 1, 5--7 and
10--13) is not used here. A date from the model means the forecast plus a margin
chosen on validation (Section~\ref{sec:stakeholders}), which the slides did not
yet have; the test is the out-of-time test window; a finding or a caveat is a
stated result or limitation; and the platform's own date is described as what it
is, a commercial choice, padded because a broken promise costs returns and
retention \citep{cui2024promise}. The verbs the slides use for what a list or a
flag does with an unhappy customer or a late parcel (slides 1, 6, 7, 9, 10 and 13) mean
only that the list contains them: precision and recall count customers who went
on to rate the order 1--2 stars and parcels that arrived late, not ratings or
deliveries that were changed. ``Pause the review email'' and ``hold the review survey until
actual arrival'' (slides 6, 10, 11 and 13) and ``Saved before 1$\star$!''
(slide 13) mean reaching the customer before they rate the order and asking for
the review once the order is complete; every customer is still asked, and no
review is filtered, delayed in publication or rewarded. ``Live validation'' and
``proves'' (slide 11) mean that the held-out days show a pattern; only the pilot
of Section~\ref{sec:pilot} would measure what a contact changes.
